\section{How it is evaluated}\label{sec:protocol}

This section gives the tasks, the references, the metrics and budgets, and the protocol, including what is held out for DFG v1.3; \cref{app:eval} gives the evaluation in full, and \cref{app:names} collects the names.

\subsection{Tasks}\label{sec:tasks}

\paragraph{Toy Boolean targets.} Inputs are uniform on $\cube$ with $d=30$, and labels are the $\pm1$ values of a target $f$. There are 11 suites, 138 targets in all, defined in \repo{study/boolfuncs.py} and listed with their status in \cref{tab:status}. The data are the networks' own (\repo{study/discrete_learner/harness.py}). For each target and seed, the training set of size $n$ is the first $n$ rows of one draw of 16384 points, so training sets are nested in $n$, and each target has one fixed test set of 4096 points. \code{designed} (21 targets) and \code{random} (20) hold the study's original juntas, parities, halfspaces, compositions and small formulas; \code{dlheld}, \code{dlheld2} and \code{dlheld3} (29) add dense halfspaces, multiplexers, formulas and symmetric functions; \code{dlsymdev}, \code{dlheld4} and \code{dlheld5} (32) symmetric and signed-sum targets. \code{dlheld6}, \code{dlfamdev} and \code{dlheld7} (12 each) are three draws from the \emph{blind family}, a random generator of signed-sum targets written in Part 9 by an agent that had not seen DFG's sum-block design (\repo{study/part9/heldout6/FAMILY.md}): nine slots of functions of a signed sum of 4--14 inputs, alone or with a parity, a gate, an inner term or a second sum, with accept sets that are thresholds, bands, intervals, residue classes mod 3 or 4, or arbitrary, and three non-symmetric controls (\cref{app:suites}). DFG v1.3 was developed on \code{dlheld6} and \code{dlfamdev} (\cref{sec:part11}), so \code{dlheld7} tests it on fresh targets from a distribution its designer had seen.

The \emph{sum-structured targets} are the 54 targets of \code{dlsymdev}, \code{dlheld4}--\code{dlheld7} and \code{dlfamdev} other than their 14 non-symmetric controls. Every suite puts a target's relevant inputs first. DFG alone and \code{lowdeg} therefore see the columns in a fresh random order in every fit, drawn from the fit's random stream (\code{harness.Permuted}), as v2 and v3 did, so that tie-breaking by column index cannot help them.

\paragraph{Input dimension.} Q4 uses the \code{designed} targets on which the networks were trained at other input dimensions in Part 3: \code{maj9}, \code{majpar}, \code{mux}, \code{par2}, \code{par3}, \code{par4} and \code{stair5} at $d=15$ and $d=60$, and \code{maj9}, \code{mux}, \code{par2} and \code{par3} at $d=120$ (7, 7 and 4 targets). Each is the same function of the first inputs as at $d=30$, and its data are drawn in the same way.

\paragraph{Tabular data.} DFG runs on the two of the study's 26 PMLB datasets \cite{olson2017pmlb} that have only binary or categorical columns and two classes, kr\_vs\_kp and mushroom, with the columns one-hot coded as $\pm1$ (\cref{app:tabular}). It runs neither on images nor on real-valued columns.

\subsection{References}\label{sec:refs}

\paragraph{Networks on the toy targets.} These are the companion paper's $\mu$P two-hidden-layer ReLU MLPs, $d$-128-128-1 \cite{gdcircuit2026,yang2021feature}: Adam \cite{kingma2015adam} at learning rate $10^{-3}$, batches of 64 sampled from the fixed training set, 20000 steps (40000 at $d=120$), read at the final step. There are five variants: \code{nets} (square loss, the only variant run at $d\ne30$), \code{nets_wd} (square loss with decoupled weight decay 0.1 \cite{loshchilov2019adamw}), and \code{nets_cewd}, \code{nets_cewd0.3} and \code{nets_cewd1.0} (logistic loss with weight decay 0.1, 0.3 and 1.0), run on different suites (\cref{app:refs}); \code{dlfamdev} has none. \code{nets_best} is, per target, the variant with the smallest $\nstar$ among the 2 to 5 run on its suite (on \code{dlheld7}, the 3 that Part 10 pre-registered): a hindsight choice that favours the networks. The networks are one recipe, not tuned per target or suite, so the comparison is with a standard gradient-trained network, not with the best network for each target.

\paragraph{No-search baselines.} \code{lin} is L2-regularized logistic regression on the raw inputs, its penalty chosen by cross-validation, and \code{hgb} scikit-learn's histogram gradient boosting at default settings \cite{pedregosa2011sklearn} (both in \repo{study/hybrid/refs.py}; details in \cref{app:refs}). Neither searches an explicit feature dictionary: \code{lin} searches no features at all, and \code{hgb} only through greedy splits on single inputs. So they show which targets need no dictionary search. Both were run in Parts 8--10, with the columns in their natural order; \code{hgb} has no runs on \code{designed}, \code{dlsymdev} and \code{dlfamdev}, and \code{lin} none on \code{dlfamdev}.

\paragraph{The low-degree algorithm: \code{lowdeg}.} \code{lowdeg} (\repo{study/discrete_learner/sparse_fourier.py}) is Part 6's first feasibility check, run unchanged with the inputs permuted per fit: the low-degree algorithm \cite{linial1993constant} with a significance threshold, which keeps every character of degree $\le4$ whose empirical coefficient on the labels passes a union bound over its degree, tests products of kept characters in the same way, and predicts with the sign of the kept expansion (\cref{app:refs}). It shares DFG's global characters up to degree 4 and the form of its test, and nothing else: no signed-sum blocks, no residual rounds and no refitted readout.

\paragraph{Portfolios that contain DFG.} In v3 (\code{hyb_portfolio3}, Part 10 \cite{portfoliolearner2026}), DFG v1.3 ran as two members, its default and (at $n\le512$) its forced variant, next to LightGBM, ExtraTrees, $k$-nearest-neighbour votes and logistic regression, with a share of the budget, and a convex stacker fitted on out-of-fold predictions weighted them (\repo{study/part10/hyb_portfolio3/learner.py}; \cref{app:v3}). Alone, DFG chooses between its two variants itself, adds its own kernel background and has the whole budget. v3's code ran on every toy suite with the inputs permuted per fit, so DFG alone and v3 can be compared target by target, which measures what the rest of the portfolio adds on Boolean inputs. v2 and v1 (Parts 9 and 8) are the earlier portfolios, with DFG v1.2 and v1.1 \cite{discretelearners2026}; earlier versions of DFG and the tabular references are in \cref{app:refs,app:tabular}.

\subsection{Metrics and budgets}\label{sec:metrics}

\paragraph{Toy targets.} $\nstar$ is the training-set size at which the seed-mean test error reaches 0.1, interpolated in $\log n$ between the grid points $n=16,32,\dots,16384$. It is ``never'' if the error does not reach 0.1 by $n=16384$. DFG alone runs on the networks' seeds and training sets: 3 seeds per target, 2 on \code{random}.
\begin{itemize}
\item A comparison of learners $A$ and $B$ on a target is the ratio $\log_2(\nstar_A/\nstar_B)$, negative when $A$ needs less data.
\item $A$ is \emph{clearly better} when the ratio is below $-0.25$ or only $A$ reaches 10\% error, and \emph{clearly worse} in the mirror cases.
\item A suite mean is the mean ratio over the targets both learn. Pooled summaries pool the per-target ratios over suites, rather than averaging suite means.
\end{itemize}
The verdict has no noise term: $\nstar$ is read from test errors averaged over 2 or 3 training sets per $(\text{target},n)$, with no confidence interval, and the $\pm0.25$ band is a convention, not a significance test. It is smaller than the spread between single seeds: DFG's $\nstar$ on two different single training sets differ by a mean $|\log_2$ ratio$|$ of 0.339 (23 of 65 targets beyond $\pm0.25$), refits with another column order by 0.012 (\cref{sec:abl-single}). Per-target verdicts rest on the 2--3-seed means and are not significance calls. These are the definitions of Parts 8--10 (\repo{study/hybrid/compare_hybrid.py}), and \repo{study/part11/compare_dfg.py} computes every comparison from the result files.

\paragraph{Input dimension.} At $d=15$, 60 and 120, DFG and the networks have seeds 0 and 1, and at $d=30$ both have 3; at $d=60$ and 120 the networks' $\nstar$ is read on their own grids, which reach up to 65536 points. The growth of $\nstar$ with $d$ is summarized by the least-squares slope of $\log\nstar$ on $\log d$, fitted for each learner over the values of $d$ at which both reach 10\% error.

\paragraph{Budgets.} Compute is counted in FLOPs, a multiply-add counting 2; everything a fit does is included, DFG's cross-validation fits too, and prediction is excluded. The unit is one training run of one toy network: $6P\cdot64\cdot T+10P\cdot T$ FLOPs for $P=128d+16768$ parameters and $T$ steps (\code{harness.gd_flops}), $1.62\times10^{11}$ at $d=30$ with $T=20000$; $1.47$ and $1.93\times10^{11}$ at $d=15$ and 60; and $5.06\times10^{11}$ at $d=120$, where the networks ran 40000 steps. Each network variant spends exactly one unit per fit, and \code{nets_best} is a hindsight minimum over up to five such runs. DFG spends what its own budget allows (\cref{sec:smalln}); the other references run their standard recipes uncapped. CPU time is the process time of a fit, rough because it was measured on a shared machine (\cref{app:compute}); the networks' result files hold no timing, so the compute comparison with them is in FLOPs only.

\subsection{The Part 11 protocol, and what is held out}\label{sec:part11}

\paragraph{Pre-registration.} The Part 11 protocol, \repo{study/part11/PROTOCOL.md}, was written and committed (\texttt{aca1f3c8}) together with the runners before any Part 11 result was produced, except one disclosed timing check, and its pre-registered part has not changed since (commits and times in \cref{app:runs}). It fixed the object and how it is run, seven questions with their runs, suites and seeds, the rules of Parts 8--10, a rule for crashes and what is held out (\cref{app:questions}). One post-hoc addendum, which defines the core run of \cref{sec:core}, was added after every pre-registered run had finished.

\paragraph{The object.} DFG v1.3 is the file \repo{study/part10/hyb_portfolio3/dfg13.py}, frozen in Part 10 with md5 \texttt{69cabe5c50345cb24079198ed52b562f}. Part 11 runs it alone, inside \code{harness.Permuted}, as the run \code{dfg13-perm}: \code{DFG(readout="logit").fit(X, y, rng)}, with every other constructor argument at its default; the runners \repo{study/part11/eval_dfg.py} and \repo{study/part11/dfg_real.py} check the md5. No line of it changes, and nothing is tuned. In the rest of the paper ``DFG'' means DFG v1.3 run alone, and ``DFG alone'' is written where the contrast with v3's DFG member matters; the reasons for choosing this object are in \cref{sec:intro}.

\paragraph{Questions and runs.} The protocol fixed seven questions: Q1, \code{dlheld7}; Q2, every suite, v3 included; Q3, \code{lowdeg}; Q4, the input dimension; Q5, 12 ablations; Q6, tabular data; and Q7, descriptive analyses of what DFG learns, checked by independent verifiers. \Cref{sec:results} reports Q1--Q4 and Q6, \cref{sec:ablations} Q5, and \cref{sec:interp} Q7; \cref{app:questions} lists each question's runs, suites and seeds.

\paragraph{What is held out.} For DFG v1.3, only \code{dlheld7} is held out (\cref{tab:status}). It was generated in Part 10 from the blind family and fixed in a commit (\texttt{dc70266c}) before any Part 10 design work; Part 10's builders, the agents that designed its components, were not allowed to load, generate or inspect it, the one that built v1.3 recorded that it saw only the one-line signature of the suite's generator function, and no result of a Part 10 learner on it was read before v3 was frozen \cite{portfoliolearner2026}.
\begin{itemize}
\item \emph{Everything else is in-sample.} Every other suite was used in developing some version of DFG, or was held out only for an earlier version, and several of DFG's parts were added after failures on, or with knowledge of, such targets: the degree-5/6 search and the spectral block (\code{dlheld}'s \code{par5} and majxor), the logistic readout (\code{thr30}), sum blocks (\code{sym9}, \code{dlheld3}) and weighted blocks (h6\_10, \code{dlheld6}). Those suites still measure the learner, but they do not test it, and some selection in its favour is expected there \status{Heuristic}. The Q4 targets are \code{designed} targets, and kr\_vs\_kp is a development set; mushroom was not used to develop any version of DFG, but v1--v3 ran DFG on it, and their results were known. The ablations and the analyses (Q5, Q7) describe the frozen learner; they are not tests of a design.
\item \emph{The object was chosen post hoc.} It was chosen knowing v3's results on every suite, including \code{dlheld7}, where DFG is v3's toy member. DFG alone's results there were new, apart from a timing check on two targets, but they are not independent of v3's, which were known. The choice is post hoc; the code is not. The timing check and the other disclosures are in \cref{app:disclosures}.
\end{itemize}
